% GENERATED FILE. Edit canonical lessons, then run npm run generate:books.
% canonical-source-hash: fnv1a64:29aa6bc4d6d387d0
% canonical-lessons: FR-C300-secours, FR-C300-secourir, FR-C300-reanimer, FR-C300-garrot, FR-C300-civiere

\chapter{Les secours}
\label{ch:fr-les-secours}
\hlchaptermodality{300}

\begin{chapteropening}
\textbf{By the end of this chapter:} \emph{I can describe a rescue: the emergency services, coming to help, reviving someone, a tourniquet and a stretcher.}

Complete the last lesson of chapter 300: I can describe a rescue: the emergency services, coming to help, reviving someone, a tourniquet and a stretcher.
\end{chapteropening}

\section[les secours]{les secours — the emergency services, help}

\label{lesson:FR-C300-secours}



\begin{quote}
Before the new one: say the French for the haematoma, then the French for the wound.
\end{quote}

\subsection*{You'll want to know: les secours}
\textbf{les secours} — \textquotedblleft{}the emergency services, help\textquotedblright{}.

\begin{quote}
Les secours sont arrivés en dix minutes.
\end{quote}

\begin{quote}
\emph{The emergency services arrived in ten minutes.}
\end{quote}

\emph{Au secours !} — help! \emph{Les premiers secours} — first aid.

\begin{cousinweb}
\emph{secours} is inherited from Latin \emph{succursus} \textquotedblleft{}help\textquotedblright{} (EtymDB), from \emph{succurrere} \textquotedblleft{}to run to help\textquotedblright{}.
\end{cousinweb}

\begin{grammarlens}[title={What you've built}]
One more word for telling, explaining and describing at B1.
\end{grammarlens}

\subsection*{Your turn}
\begin{itemize}
\raggedright
  \item \emph{Say it:} \emph{les secours}
  \item \emph{Say it:} \emph{les secours}, once more
  \item \emph{Say it:} a sentence of your own with \emph{les secours}
  \item \emph{Recall:} say the French for the haematoma, then the French for the wound, then say \emph{les secours} again
\end{itemize}

\begin{tcolorbox}[breakable,colback=teal!4,colframe=teal!35!black,title={Before you move on}]
What does les secours mean? (\textquotedblleft{}The emergency services, help\textquotedblright{}.) Say it once more.
\end{tcolorbox}
\section[secourir]{secourir — to rescue, to help}

\label{lesson:FR-C300-secourir}



\begin{quote}
Before the new one: say the French for the wound, then the French for the emergency services.
\end{quote}

\subsection*{You'll want to know: secourir}
\textbf{secourir} — \textquotedblleft{}to rescue, to help\textquotedblright{}.

\begin{quote}
Des voisins ont secouru la vieille dame.
\end{quote}

\begin{quote}
\emph{Neighbours helped the old lady.}
\end{quote}

Conjugated like \emph{courir}: \emph{il a secouru}.

\begin{cousinweb}
\emph{secourir} is inherited from Latin \emph{succurrere} \textquotedblleft{}to run to help\textquotedblright{}, from \emph{sub} and \emph{currere} \textquotedblleft{}to run\textquotedblright{} (Lewis \& Short).
\end{cousinweb}

\begin{grammarlens}[title={What you've built}]
One more word for telling, explaining and describing at B1.
\end{grammarlens}

\subsection*{Your turn}
\begin{itemize}
\raggedright
  \item \emph{Say it:} \emph{secourir}
  \item \emph{Say it:} \emph{secourir}, once more
  \item \emph{Say it:} a sentence of your own with \emph{secourir}
  \item \emph{Recall:} say the French for the wound, then the French for the emergency services, then say \emph{secourir} again
\end{itemize}

\begin{tcolorbox}[breakable,colback=teal!4,colframe=teal!35!black,title={Before you move on}]
What does secourir mean? (\textquotedblleft{}To rescue, to help\textquotedblright{}.) Say it once more.
\end{tcolorbox}
\section[réanimer]{réanimer — to resuscitate, to revive}

\label{lesson:FR-C300-reanimer}



\begin{quote}
Before the new one: say the French for the emergency services, then the French for to rescue.
\end{quote}

\subsection*{You'll want to know: réanimer}
\textbf{réanimer} — \textquotedblleft{}to resuscitate, to revive\textquotedblright{}.

\begin{quote}
Les pompiers ont réanimé l'homme.
\end{quote}

\begin{quote}
\emph{The firefighters revived the man.}
\end{quote}

\emph{Le service de réanimation} — intensive care.

\begin{cousinweb}
\emph{réanimer} is built on \emph{animer}, a learned borrowing from Latin \emph{animāre} \textquotedblleft{}to give life\textquotedblright{}, from \emph{anima} \textquotedblleft{}breath, life\textquotedblright{} (Lewis \& Short).
\end{cousinweb}

\begin{grammarlens}[title={What you've built}]
One more word for telling, explaining and describing at B1.
\end{grammarlens}

\subsection*{Your turn}
\begin{itemize}
\raggedright
  \item \emph{Say it:} \emph{réanimer}
  \item \emph{Say it:} \emph{réanimer}, once more
  \item \emph{Say it:} a sentence of your own with \emph{réanimer}
  \item \emph{Recall:} say the French for the emergency services, then the French for to rescue, then say \emph{réanimer} again
\end{itemize}

\begin{tcolorbox}[breakable,colback=teal!4,colframe=teal!35!black,title={Before you move on}]
What does réanimer mean? (\textquotedblleft{}To resuscitate, to revive\textquotedblright{}.) Say it once more.
\end{tcolorbox}
\section[le garrot]{le garrot — the tourniquet}

\label{lesson:FR-C300-garrot}



\begin{quote}
Before the new one: say the French for to rescue, then the French for to resuscitate.
\end{quote}

\subsection*{You'll want to know: le garrot}
\textbf{le garrot} — \textquotedblleft{}the tourniquet\textquotedblright{}.

\begin{quote}
Un médecin a posé un garrot sur sa jambe.
\end{quote}

\begin{quote}
\emph{A doctor put a tourniquet on his leg.}
\end{quote}

Also the withers of a horse.

\begin{cousinweb}
The origin of \emph{garrot} is not settled; it is sometimes linked to Provençal \emph{garra} \textquotedblleft{}leg\textquotedblright{} (EtymDB).
\end{cousinweb}

\begin{grammarlens}[title={What you've built}]
One more word for telling, explaining and describing at B1.
\end{grammarlens}

\subsection*{Your turn}
\begin{itemize}
\raggedright
  \item \emph{Say it:} \emph{le garrot}
  \item \emph{Say it:} \emph{le garrot}, once more
  \item \emph{Say it:} a sentence of your own with \emph{le garrot}
  \item \emph{Recall:} say the French for to rescue, then the French for to resuscitate, then say \emph{le garrot} again
\end{itemize}

\begin{tcolorbox}[breakable,colback=teal!4,colframe=teal!35!black,title={Before you move on}]
What does le garrot mean? (\textquotedblleft{}The tourniquet\textquotedblright{}.) Say it once more.
\end{tcolorbox}
\section[la civière]{la civière — the stretcher}

\label{lesson:FR-C300-civiere}



\begin{quote}
Before the new one: say the French for to resuscitate, then the French for the tourniquet.
\end{quote}

\subsection*{You'll want to know: la civière}
\textbf{la civière} — \textquotedblleft{}the stretcher\textquotedblright{}.

\begin{quote}
Ils l'ont transporté sur une civière.
\end{quote}

\begin{quote}
\emph{They carried him on a stretcher.}
\end{quote}

Carried by two people.

\begin{cousinweb}
The origin of \emph{civière} is not settled (EtymDB lists no parent).
\end{cousinweb}

\begin{grammarlens}[title={What you've built}]
One more word for telling, explaining and describing at B1.
\end{grammarlens}

\subsection*{Your turn}
\begin{itemize}
\raggedright
  \item \emph{Say it:} \emph{la civière}
  \item \emph{Say it:} \emph{la civière}, once more
  \item \emph{Say it:} a sentence of your own with \emph{la civière}
  \item \emph{Recall:} say the French for to resuscitate, then the French for the tourniquet, then say \emph{la civière} again
\end{itemize}

\begin{tcolorbox}[breakable,colback=teal!4,colframe=teal!35!black,title={Before you move on}]
What does la civière mean? (\textquotedblleft{}The stretcher\textquotedblright{}.) Say it once more.
\end{tcolorbox}
